% Planning van één labo. Gebruikt door de syllabus van dit labo en door de syllabus voor labo 3-5.
\textbf{Doel:} een dataset inladen en verkennen, rijen en kolommen
selecteren en filteren, nieuwe kolommen toevoegen en de data opkuisen.

\begin{schedule}
13:00--13:20 & Opwarmer: waarom niet gewoon Excel? & De ruwe tekst van een CSV-bestand bekijken en zien wat een tabelbibliotheek toevoegt. \\
13:20--14:00 & Inladen en verkennen & \code{read\_csv}, \code{.head()}, \code{.shape}, \code{.dtypes}, \code{.info()}, \code{.describe()}, \code{mo.ui.table}. \\
14:00--14:50 & Selecteren en filteren & Kolommen selecteren, \code{.loc} met logische voorwaarden, voorwaarden combineren, \code{sort\_values}. \\
14:50--15:05 & Pauze & \\
15:05--15:50 & Nieuwe kolommen & Tekst naar getal, leeftijd uit een geboortedatum, snelheid uit een sprinttijd, BMI. \\
15:50--17:00 & Opkuisen & Ontbrekende waarden (\code{dropna} of \code{fillna}), tekstlabels, eenheden, dubbels, één volledige opkuis-pipeline. \\
17:00--17:30 & Uitdaging & De 5 snelste atleten per positie. Exit ticket. \\
\end{schedule}

\textbf{Oefeningen:} (1) \code{tests.csv} inladen; (2) aanvallers jonger
dan 16; (3) verdachte sprinttijden; (4) BMI; (5) leeftijd in jaren;
(6) de posities opkuisen; (7) de eenheden van de sprinttijden
corrigeren; (8) de dubbele atleet verwijderen; \textbf{uitdaging:} de
5 snelsten per positie.
